\chapter{Conclusion}
\label{ch:conclusion}

This thesis set out to realize the promise of energy-based controllable text generation, in which a frozen autoregressive language model is treated as an energy function and steered at inference time by sampling from it with gradient-guided Markov chain Monte Carlo. Across a study of one hundred and forty-five sampling configurations, five energy functions, and two samplers, it found that the central premise of the framework does not hold: the frozen likelihood of an autoregressive language model is not a usable energy function for sampling on discrete text.

The contribution of the thesis is not this negative observation alone but the mechanistic account that accompanies it, established by independent means and supported at every step by direct measurement. Following the true gradient was shown to be statistically indistinguishable from following a random direction of equal magnitude, on all five energy functions, once a normalization confound was removed by an ablation designed for the purpose. This gradient fallacy was explained by a linearization failure: the first-order surrogate that the discrete sampler uses to rank candidate tokens is uncorrelated with the true change in energy, because the smallest admissible discrete step is far larger than the radius over which the local linear model is valid. The Metropolis--Hastings correction, applied faithfully, was shown to paralyze the continuous sampler, and the trace analysis located the cause precisely in the proposal term, which collapses because the nearest-neighbour projection makes the drift discontinuous at every cell boundary. A likelihood trap was measured on the study's own models, showing that the lowest-energy text is the most degenerate, so that success at the optimization objective is failure at the underlying goal. Turning to amortized inference, the GFlowNet experiments exhibited three distinct failure modes, and a unifying experiment showed that Langevin sampling on the GFlowNet-tuned energy is indistinguishable from sampling on the untuned base, demonstrating that the pathology survives amortization and is therefore a property of the energy rather than of any sampler. A constrained-generation extension confirmed that the constraint gradient carries no reliable steering signal on this landscape.

Taken together, these findings locate the cause of the failure in the training objective of autoregressive language models, which shapes accurate local conditional distributions and never a smooth global energy over sequences. This diagnosis makes a falsifiable prediction, that a diffusion language model whose objective does shape a score function should not fail in the same way, and the thesis sets out that positive-control experiment as the natural next step. What the thesis provides to the larger effort of building inference-time controllable generation on frozen models is therefore a precise characterization of where and why the most common approach breaks down, together with an experiment capable of confirming or refuting the explanation. A negative result about a method, when it is accompanied by a mechanism, a measurement, and a falsifiable hypothesis, is a positive result about understanding, and it is in that sense that the work is offered.
